\documentclass{article}
\usepackage[utf8]{inputenc}
\usepackage[normalem]{ulem} 
\usepackage[dvipsnames]{xcolor}
\definecolor{leanblue}{RGB}{28, 54, 86}
\newcommand{\comment}[1]{}
\usepackage{tabularx}
\usepackage{verbatim}
\usepackage{hyperref}
\usepackage{mathtools}
\usepackage{float}
\usepackage{subcaption}
\usepackage{textcomp}
\usepackage[a4paper, margin={1in}]{geometry}
\usepackage[english]{babel}
\usepackage{graphicx}
\usepackage{amsmath}
\usepackage{amssymb}
\usepackage[bb=boondox]{mathalfa}
\usepackage{xparse}
\usepackage{fontawesome5}
\usepackage{amsfonts}
\usepackage{fontspec}
\usepackage{bbm}
\setmonofont{FreeMono}
\usepackage{minted}
\setminted{breaklines=true, breakanywhere=true}
\fvset{breaklines=true, breakanywhere=true}
\newminted[leancode]{lean4}{fontsize=\footnotesize}
\usemintedstyle{tango}

\newcolumntype{Y}{>{\centering\arraybackslash}X}

\usepackage{amsthm}
\usepackage{authblk}
\usepackage{biblatex}
\usepackage{mathpartir}
\usepackage{stmaryrd}
\usepackage{amssymb}
\newtheorem{lemma}{Lemma}
\newtheorem{theorem}{Theorem}
\newtheorem{definition}{Definition}
\newtheorem{example}{Example}
\addbibresource{./refs.bib}

\newcommand{\dcc}{\(\text{DCC}\)}
\newcommand{\dccbf}{\(\textbf{DCC}\)}
\newcommand{\dcce}{\text{DCC}}
\newcommand{\dccebf}{\textbf{DCC}}
\newcommand{\ccc}{\(\text{CCC}\)} % Cartesian-closed category
\newcommand{\cccbf}{\(\textbf{CCC}\)} % Cartesian-closed category
\newcommand{\ccl}{\(\text{CCL}\)} % Categorical combinatory logic
\newcommand{\cclbf}{\(\textbf{CCL}\)} % Categorical combinatory logic
\newcommand{\ccle}{\text{CCL}}
\newcommand{\cclebf}{\textbf{CCL}}
\newcommand{\cclk}{\(\text{CCL}_\text{K}\)} % Simply-typed categorical combinatory logic
\newcommand{\cclkbf}{\(\textbf{CCL}_\textbf{K}\)} % Simply-typed categorical combinatory logic
\newcommand{\cclke}{\text{CCL}_\text{K}}
\newcommand{\sk}{\(\text{SK}\)}
\newcommand{\skbf}{\(\textbf{SK}\)}
\newcommand{\ske}{\text{SK}}
\newcommand{\skebf}{\textbf{SK}}
\newcommand{\cclkebf}{\textbf{CCL}_\textbf{K}}
\newcommand{\llb}{\llbracket}
\NewDocumentCommand{\comb}{o o o o o}{
  \IfNoValueTF{#5}{
    \IfNoValueTF{#4}{
      \IfNoValueTF{#3}{(#1\ #2)}{(#1\ #2\ #3)}}{(#1\ #2\ #3\ #4)}}{(#1\ #2\ #3\ #4\ #5)}}
\newcommand{\rrb}{\rrbracket}
\newcommand{\leanlink}[2]{
  \hfill \href{#1}{\textcolor{leanblue}\faCheckSquare~[Formalized: \texttt{#2}]}%
}
\newcommand{\la}{\langle}
\newcommand{\ra}{\rangle}
\NewDocumentCommand{\pair}{o m}{
  \IfNoValueTF{#1}{\la #2 \ra}{\llb #1 \rrb \la #2 \ra}
}
\NewDocumentCommand{\comp}{o m m}{
  \IfNoValueTF{#1}{#2 \circ #3}{
    \Ifstr{#1}{'}{
      #2 \circ' #3
    }{%
      \Ifstr{#1}{2}{#2 \circ \circ #3
      }{#2 \circ' #3}
    }
  }
}
\NewDocumentCommand{\cc}{o m}{
  \IfNoValueTF{#1}{\llb #2 \rrb}{\llb #2 \rrb_{#1}}
}
\NewDocumentCommand{\coup}{o m}{
  \IfNoValueTF{#1}{(#2)
  }{\llb #1 \rrb (#2)}
}
\newcommand{\arr}{\rightarrow}
\newcommand{\sort}[1]{\text{Sort}\ #1}
\newcommand{\type}[1]{\text{Type}_{#1}}
\newcommand{\barr}{\arr_{\beta}}
\newcommand{\barrc}[1]{\arr_{\beta^{#1}}}
\newcommand{\barrn}{\arr_{\beta}^{!}}
\newcommand{\earr}{\arr_{\eta}}
\newcommand{\affarr}{\rightarrowtail}
\newcommand{\affetaarr}{\rightarrowtail_{\eta}}
\newcommand{\etaarr}{\rightarrow_{\eta}}
\newcommand{\dom}[1]{\text{dom}(#1)}
\newcommand{\domtwo}[1]{\text{dom}_{2}(#1)}
\newcommand{\arity}[2]{\text{arity} #1 = #2}
\newcommand{\efn}[2]{\text{F}_{\eta}\ #1\ #2}
\newcommand{\ddom}{\text{Domain}\ }
\newcommand{\dddom}{\text{Domain}_{2}\ }
\newcommand{\hole}{\text{Hole}\ }

\title{Combinator Calculus: Decidable Type Checking in Extensional Type Theory}
\author[1]{Lexi Aiello}
\author[2]{Eric Klavins}
\affil[1]{Department of Philosophy, University of Washington}
\affil[2]{Electrical \& Computer Engineering Department, University of Washington}

\begin{document}

\maketitle

\begin{abstract}
  \textit{TODO}
\end{abstract}

\section{TODO}

Should be structured around building up to the fundamental theorem: extensional equality reflection.

\begin{itemize}
  \item \sout{Update results section}
   \begin{itemize}
      \item \sout{It should be emphasized that equality reflection and other metatheorems (e.g., subject-reduction, progress) are computable definitions, rather than theorems in \(\textbf{Prop}\).}
  \end{itemize}
  \item \sout{Syntax}
    \begin{itemize}
      \item \sout{Give a presentation in BNF notation}
    \end{itemize}
  \item Conversion rules
  \begin{itemize}
    \item \sout{Emphasize that we use call by value operational semantics as usual}
    \item \sout{Explain how \(Eval\) is essentially the recursor for \(\Sigma\)}
    \item \sout{Congruences}
    \begin{itemize}
      \item \sout{Explain combinator points, maybe do this in syntax section though}
    \end{itemize}
    \item \textbf{Proof}: Termination by decreasing syntactic size
    \begin{itemize}
      \item Sizes should be somewhat justified, rather than the random values found in the formalization
    \end{itemize}
    \item \textbf{Proof}: Local confluence
    \item \textbf{Proof}: Confluence by newmann's lemma
  \end{itemize}
  \item Inference rules
    \begin{itemize}
      \item Explain how \(\Pi\) indexes an environment, similarly to \(\Lambda\) among others
      \item Must explain arrows, different kinds of arrows, inductives used for representing arrow types. \textit{This can be done through a rule in the main inference rules diagram}.
      \item Inference rules should be given first as a full page diagram, and then justified for some in particular (\textit{which ones?})
    \end{itemize}
  \item Extensional equality reflection proof - \textit{note: this requires subject reduction and progress}
  \begin{itemize}
    \item \textbf{Proof}: Decidability
    \item \textbf{Proof}: Subject-reduction lemma
    \item \textbf{Proof}: Progress theorem
      \begin{itemize}
        \item Canonical values should be justified
        \item Emphasize that the rule for \(\bot{}.Rec\) is not for \(\bot\), but partial application of its recursor
        \item \textbf{Proof}: Inversion lemmas
        \item \textbf{Proof}: \(\bot\) is not inhabited by any values
       \end{itemize}
     \item \textbf{Proof}: Equality reflection for values
     \item \textbf{Proof}: Equality reflection by progress + preservation
  \end{itemize}
  \item Justification of affine combinator rules
  \item Some examples
  \item Compilation from Lean 4 (calculus of constructions subset) to DCC
  \item Comparison to SK combinators and prior approaches
  \item \textit{Last}: conclusion, somewhere in there - In this paper, we present DCC: a dependent combinatory type theory based on affine logic. We demonstrate that it is provably a counter-example to the well-known fact that type-checking is undecidable in extensional type theory (\textit{how well-known?}).
\end{itemize}

\section{Introduction}

\textit{TODO later after main sections done}

\subsection{Results}

\begin{itemize}
  \item In \ref{dccintro}, we present the dependent combinator calculus (\dccbf): our extension of the categorical combinators \autocite{cmachine} \autocite{CURIEN1986188} to extensional dependent type theory and affine logic. We have no variables and no contexts. Types are terms in DCC, as in the Calculus of Constructions with respect to the \(\lambda\)-cube \autocite{Barendregt_1991}. We describe our \(\beta\)-reduction rules and inference rules.
  \item In \(\ldots\) (\textit{todo}), we prove some basic properties of conversion in DCC: strong normalization, confluence, and uniqueness of normal forms.
  \item In \ref{decidable}, we prove that type checking is decidable in DCC.
  \item In \(\ldots\) (\textit{todo}), we prove that subject-reduction and progress hold in DCC, and that DCC is consistent.
  \item In \ref{eqreflect}, we prove that extensional equality reflection holds in DCC. The proof is a computable definition, rather than a theorem in \(\textbf{Prop}\).
  \item In \ref{compilelean}, we present a procedure for compiling the subset of Lean 4 that maps to the base Calculus of Constructions to DCC.
\end{itemize}

All of our results have been formalized in the Lean 4 proof assistant. To this date, DCC is the only calculus with extensional equality reflection and decidable type-checking. However, due to our choice of affine logic and categorical combinators, we inherit these features by construction.

\subsection{Background}

\section{Term Syntax}

The base objects are a modified superset of the base terms in \cclk{}. Products are replaced by \(\Sigma\), arrows by \(\Pi\), sorts are added, and bottom, top, and boolean types are added with the corresponding introducing terms. The syntax is given in S-expressions for clarity. A term (\(E\)) is either a constant of term arity zero (\(E_c\)), a parenthesized combinator (\(C\)) juxtaposed with its points, or a parenthesized application of two terms.

\begin{definition}\label{def:constants}
  Constants (\(E_c\)) are sort levels (written as \(\sort{u}\)), the bottom (\(\mathbb{0}\)), top (\(\mathbb{1}\)) type and its inhabitant (\(\ast\)), and the boolean type (\(\mathbb{2}\)) along with its inhabitants \(\mathbb{1}_{\mathbb{2}}\), and \(\mathbb{0}_{\mathbb{2}}\):
  \[
    \begin{gathered}
      E_c \coloneq \{\sort{u}, \mathbb{0}, \mathbb{1}, \mathbb{2}, \ast, \mathbb{1}_{\mathbb{2}}, \mathbb{0}_{\mathbb{2}}\}
    \end{gathered}
  \]
  where \(u \in \mathbb{N}\).
\end{definition}\label{def:constants}

\begin{definition}\label{def:combinator}
  Combinators are identified by the set of names \(C\):
  \[
    \begin{gathered}
      C \coloneq \{\mathbb{0}.Rec, \mathbb{2}.Rec, !<>, \$<>, *, !, \&, =, Eval, Eql, =', refl, `, Id, Fst, Snd, Swap, \Sigma, \Sigma', \Lambda, !\Lambda, \Pi, \circ, \langle \, \rangle, ,\}
      \end{gathered}
  \]
\end{definition}

\begin{definition}\label{def:syntax}
  A term (\(E\)) is either a constant \(E_c\):
  \[E_c \subset E\]
  or a parenthesized combinator \(C\) juxtaposed by its points which are also terms:
  \[\{ (c x_1 \ldots x_n) | c \in C, \{ x | x \in x_0 \ldots x_n\} \subset E \}\ \subset E\]
  or a parenthesized application of a list of juxtaposed terms:
  \[\{(f x_1 \ldots x_n) | f \in E, \{x | x \in x_0 \ldots x_n\} \subset E\} \subset E\]
\end{definition}

\begin{example}
  If \(A \in E\), then \(\comb[Id][A] \in E\).
\end{example}

\begin{example}
  If \(A, x \in E\), then \((\comb[Id][A] x) \in E\).
\end{example}

\begin{example}
  If \(u \in \mathbb{N}\), then \(\sort{u} \in E\).
\end{example}

\begin{example}
  If \(B, x, y \in E\), then \(\comb[,B][x][y] \in E\).
\end{example}

\section{Conversion Rules}

Conversion uses call-by-value operational semantics and is a modified superset of those in \ccl{}. Terms are reduced in any particular order of: head reductions, points of the term if it is a combinator, or the constituent function and argument if it is an application. We have global confluence, strong normalization, and uniqueness of normal forms.

\begin{definition}
  The identity, quotation, and composition are the same as in \cclk{}:
  
  \[
    \begin{gathered}
  \inferrule
  { }
  { (\comb[Id][A] x) \barr x} \qquad
  \inferrule
  { }
  { (\comb[`B][x] y) \barr x}\qquad
  \inferrule
  { }
  { (\comb[\circ][x][y] z) \barr (x (y z))}
  \end{gathered}
  \]
\end{definition}

\begin{definition}
  Projection over couples and currying are the same as in \ccl{}, but with type arguments added:

  \[
  \inferrule
  { }
  { (\comb[Fst][B] \comb[,B_2][x][y]) \barr x}\qquad
  \inferrule
  { }
  { (\comb[Snd][B] \comb[,B_2][x][y]) \barr y}\qquad
  \inferrule
  { }
  { (\comb[\Lambda][A][B][x] y z) \barr x \comb[,B][y][z]}
  \]
\end{definition}

\begin{definition}
  Evaluation of couples is given as the usual recursor for \(\Sigma\):

  \[
    \inferrule
    { }
    { (\comb[Eval][A][B][x] \comb[,B_2][y][z]) \barr (x y z)}
  \]
\end{definition}

\begin{definition}
  Exchange is permitted on nondependent couples as an optimization, yielding a couple quoting the first component:
  \[
    \inferrule
    { }
    { (\comb[Swap][A][B] \comb[,B_2][x][y]) \barr \comb[\comb[`B][A]][y][x]}
  \]
\end{definition}

\begin{definition}
  Boxed terms may be unboxed, mapped into banged types, and mapped over into other boxed terms:
  \[
    \inferrule
    { }
    { (\comb[*][A] \comb[\&][x]) \barr x}\qquad
    \inferrule
    { }
    { (\comb[!<>][x] \comb[\&][y]) \barr \comb[!][(x y)]}\qquad
    \inferrule
    { }
    { (\comb[\$<>][x] \comb[\&][y]) \barr \comb[\&][(x y)]}
  \]
\end{definition}

\begin{definition}
  Boxed terms can also be curried:
  \[
    \inferrule
    { }
    { (\comb[!\Lambda][A][B][x] \comb[\&][y] \comb[\&][z]) \barr (x \comb[\&][\comb[,B][y][z]])}
  \]
\end{definition}

\begin{definition}
  As opposed to its usual behavior in \ccl{}, the pair rule consumes boxed terms and duplicates its argument so as to be fully dependent:
  \[
    \inferrule
    { }
    { (\comb[\langle , \rangle B][x][y] \comb[\&][x]) \barr \comb[,(B x)][(M x)][(N x (M x))]}
  \]
\end{definition}

\begin{definition}
  Equality types can also be formed by the combinators \(='\) and \(Eql\), which index an environment determining the arguments' types:
  \[
    \inferrule
    { }
    {(\comb[='][B] x y z) \barr \comb[=(B x)][y][z]}\qquad
    \inferrule
    { }
    {(\comb[Eql][B] \comb[\&][x] \comb[\&][y]) \barr \comb[=(B x)][y][y]}
  \]
\end{definition}

\begin{definition}
  \(\Sigma\)-types can be formed from boxes:
  \[
    \inferrule
    { }
    {(\comb[\Sigma'][B][C] \comb[\&][x]) \barr \comb[\Sigma][(B x)][(C x)]}
  \]
\end{definition}

\begin{definition}
  We have a recursor for \(\mathbb{2}\):
  \[
    \inferrule
    { }
    { (\comb[\mathbb{2}.Rec][motive][x][y] z \mathbb{1}_\mathbb{2}) \barr (x z \mathbb{1}_\mathbb{2})}\qquad
    \inferrule
    { }
    { (\comb[\mathbb{2}.Rec][motive][x][y] z \mathbb{0}_\mathbb{2}) \barr (y z \mathbb{0}_\mathbb{2})}
  \]
\end{definition}

\begin{definition}
  Finally, there are pointwise congruences for combinators and congruences for application:
  \[
    \inferrule
    { c \in C \qquad i < n \qquad x_i \barr x'_i }
    { (c x_0 \ldots x_n) \barr (c x'_0 \ldots x'_n)}\qquad
    \inferrule
    { f \barr f'}
    { (f x) \barr (f' x)}\qquad
    \inferrule
    { x \barr x'}
    { (f x) \barr (f x')}
  \]
\end{definition}

Operations are added to curry and transform affine boxed terms. The pair rule is restricted affinely.

\section{Inference Rules}

\subsection{Type-Checking is Decidable}

\section{Computable Equality Reflection}

\section{Examples}

\section{Compiling Lean 4 to DCC}

\section{Comparison to SK Combinators}

\section{Limitations}

\section{Conclusion}

\textit{TODO later after main sections done}

\printbibliography

\end{document}
